\documentclass[%
fleqn,%
paper=a4paper,%
fontsize=10pt,%
open_bracket_pos=zenkakunibu_nibu,%
hanging_punctuation,%
]%
{jlreq}
\jlreqsetup{%
itemization_beforeafter_space=0pt,%
itemization_itemsep=0pt%
}
\makeatletter
\RequirePackage{luatexja}
\RequirePackage{luatexja-otf}
\RequirePackage{graphicx}
\RequirePackage{amsmath}
\DeclareRobustCommand{\metaphysicaicon}{\raisebox{-4.0pt}{\includegraphics[width=16pt]{metaphysicaicon.pdf}}}
\RequirePackage[normalem]{ulem}
\RequirePackage[explicit]{titlesec}
\titleformat{\section}[hang]{}{}{0pt}{\uuline{\raisebox{1pt}{\textsf{\thesection\quad #1}}}}[\vspace{0.35\baselineskip}]
\renewcommand{\thesection}{\S\,\arabic{section}}
\let\originalsection\section
\DeclareRobustCommand{\section}{\@ifstar{\@metaphysica@section@star}{\@metaphysica@section@nostar}}
\DeclareRobustCommand{\@metaphysica@section@star}[1]{\vspace{0.5\baselineskip}\originalsection{#1}\vspace*{-\baselineskip}}
\DeclareRobustCommand{\@metaphysica@section@nostar}[1]{\vspace{0.5\baselineskip}\originalsection{#1}}
\RequirePackage[%
truedimen,%
margin=30truemm,
includehead%
]{geometry}
\RequirePackage{lastpage}
\RequirePackage{fancyhdr}
\pagestyle{fancy}
\DeclareRobustCommand{\headertitle}[2][\metaphysicaicon]{%
\rhead[#2]{#1{}\quad\thepage{}/{}\pageref{LastPage}}%
\lhead[\thepage{}/{}\pageref{LastPage}\quad{}#1]{#2}%
\cfoot{}%
}
\RequirePackage{setspace}
\setstretch{1.155}
\DeclareRobustCommand{\linespace}{\@ifstar{\vspace{\baselineskip}}{\vspace{0.25\baselineskip}}}
\DeclareRobustCommand{\linesmash}{\@ifstar{\vspace{-\baselineskip}}{\vspace{-0.25\baselineskip}}}
\AtBeginDocument{%
\abovedisplayskip     =0.125\abovedisplayskip
\abovedisplayshortskip=0.125\abovedisplayshortskip
\belowdisplayskip     =0.125\belowdisplayskip
\belowdisplayshortskip=0.125\belowdisplayshortskip}
\setlength{\jot}{0pt}%
\setlength{\mathindent}{2\zw}%
\renewcommand{\floatpagefraction}{0.75}
\allowdisplaybreaks[2]
\RequirePackage[no-math]{fontspec}
\RequirePackage[no-math,deluxe,haranoaji]{luatexja-preset}
\RequirePackage{multicolpar}
\RequirePackage[style=iso]{datetime2}
\RequirePackage[unicode]{hyperref}
\RequirePackage{xparse}
\RequirePackage{dashbox}
\newcounter{psuedosectioncounter}
\setcounter{psuedosectioncounter}{1}
\newcounter{psuedocontentscounter}
\setcounter{psuedocontentscounter}{1}
\DeclareRobustCommand{\psuedosection}[3]{%
\hypertarget{#1}{\mbox{}}\begin{multicolpar}{2}%
\noindent\uuline{{\raisebox{1pt}{\textsf{\S\ \thepsuedosectioncounter\quad #2}}}}

\noindent\uuline{{\raisebox{1pt}{\textsf{\S\ \thepsuedosectioncounter\quad #3}}}}
\end{multicolpar}%
\stepcounter{psuedosectioncounter}%
\vspace{\baselineskip}%
}
\DeclareRobustCommand{\psuedocontents}[3]{%
\begin{multicolpar}{2}%
\noindent{\textsf{\hyperlink{#1}{\S\ \thepsuedocontentscounter\quad #2}}}

\noindent{\textsf{\hyperlink{#1}{\S\ \thepsuedocontentscounter\quad #3}}}\end{multicolpar}%
\stepcounter{psuedocontentscounter}%
}
\newenvironment{translateing}%
{\begin{multicolpar}{2}}
{\end{multicolpar}\vspace{\baselineskip}}
\DeclareRobustCommand{\maketitletranslating}%
{\maketitle\thispagestyle{fancy}
\vspace{\baselineskip}\begin{multicolpar}{2}
\textsf{English}

\noindent
\textsf{日本語 (Japanese)}
\end{multicolpar}\vspace{\baselineskip}}
\NewDocumentCommand\macroexplanation{v}{%
\noindent\hspace*{\fill}{\texttt{#1}}\hspace*{\fill}\linespace%
}
\NewDocumentEnvironment{macroexample}{O{0.625} +b}{%
\noindent\hspace*{\fill}\dbox{\parbox{#1\textwidth}{%
#2%
}}\hspace*{\fill}}%
{\vspace{\baselineskip}}
\NewDocumentEnvironment{macroexample*}{O{0.625} m +b}{%
\noindent\hspace*{\fill}\dbox{\parbox{#1\textwidth}{%
\vspace{-0.5\baselineskip}\begin{#2}%
#3
\end{#2}%
}}\hspace*{\fill}}
{\vspace{\baselineskip}}
\let\code\texttt
\setlength{\fboxsep}{1em}
\setstretch{1.05}
\DeclareRobustCommand{\commandtojskip}{\hspace{2.40554pt plus 1.49994pt minus 0.59998pt}}
\RequirePackage{listings, jlisting}
\lstset{
  language=[LaTeX]TeX,
  basicstyle={\ttfamily},
  identifierstyle={\small},
  commentstyle={\small\itshape},
  keywordstyle={\small\bfseries},
  ndkeywordstyle={\small},
  stringstyle={\small\ttfamily},
  frame=single,
  breaklines=true,
  columns=[l]{fullflexible},
  stepnumber=1,
  xrightmargin=0.1709\textwidth,
  xleftmargin=0.1709\textwidth,
  lineskip=-0.5ex
}
\RequirePackage{bxtexlogo}
\RequirePackage{shortvrb}
\MakeShortVerb{\|}
\RequirePackage[luacircled]{inlinelabel}
\makeatother
%
\hypersetup{%
bookmarksnumbered=true,%
colorlinks=true,%
linkcolor=blue,%
urlcolor=blue,%
setpagesize=false,%
pdftitle={The inlinelabel package},%
pdfauthor={Yukoh KUSAKABE},%
pdfsubject={The inlinelabel package},%
pdfkeywords={TeX LaTeX inlinel label}}
\title{The \code{inlinelabel} package:\\[0.25\baselineskip]
assign equation numbers to inline equations}
\author{Yukoh KUSAKABE}
%\author{Y\=uk\=o KUSAKABE}
\date{\today}
\headertitle[Yukoh KUSAKABE\quad\metaphysicaicon]{The \code{inlinelabel} package}
%\headertitle[Y\=uk\=o KUSAKABE\quad\metaphysicaicon]{The \code{inlinelabel} package}
\begin{document}
\maketitletranslating

\begin{translateing}
This package can assign equation numbers to inline equations. When Japanese is supported, you can switch to circled equation numbers.

このパッケージは，インライン数式に数式番号を振ることができます。また，日本語を扱えるときには，丸で囲まれた数式番号に切り替えることができます。
\end{translateing}

%\psuedocontents{inlinelabel}
%{Package \code{inlinelabel}}
%{\code{inlinelabel} パッケージ}

\psuedocontents{Requirements}{System Requirements}{前提条件}

\psuedocontents{Installation}{Installation}{インストール}

\psuedocontents{Loading}{Loading}{読み込み}

\psuedocontents{Usage}{Usage}{使用方法}

\psuedocontents{moreinfo}{For More Information}{問い合わせ・詳しくは}

%\psuedosection{inlinelabel}{Package \code{inlinelabel}}{\code{inlinelabel} パッケージ}
%
%\begin{translateing}
%This package can assign equation numbers to inline equations. When Japanese is supported, you can switch to circled equation numbers.
%
%このパッケージは，インライン数式に数式番号を振ることができます。また，日本語を扱えるときには，丸で囲まれた数式番号に切り替えることができます。
%\end{translateing}

\psuedosection{Requirements}{System Requirements}{前提条件}

\begin{translateing}
\textbullet\ \LaTeXe\ format\\
\textbullet\ \code{amsmath} package\\
\textbullet\ \code{refcount} package\\
\textbullet\ \pTeX\ engine (\code{[circled]} only)\\
\textbullet\ \code{japanese-otf} package (\code{[circled]} only)\\
\textbullet\ \LuaTeX\ engine (\code{[luacircled]} only)\\
\textbullet\ \scalebox{0.95}[1]{\code{luatexja-otf} package (\code{[luacircled]} only)}

\noindent
\textbullet\ \LaTeXe フォーマット\\
\textbullet\ \code{amsmath} パッケージ\\
\textbullet\ \code{refcount} パッケージ\\
\textbullet\ \pTeX エンジン（\code{[circled]}使用時）\\
\textbullet\ \scalebox{0.9}[1]{\code{japanese-otf} パッケージ（\code{[circled]}使用時）}\\
%\textbullet\ \code{luatexja-otf} パッケージ\\
%\hfill （\code{[luacircled]}使用時）
\textbullet\ \LuaTeX エンジン（\code{[luacircled]}使用時）\!\\
\textbullet\ \scalebox{0.8375}[1]{\code{luatexja-otf} パッケージ（\code{[luacircled]}使用時）}
\end{translateing}

\newpage
\psuedosection{Installation}{Installation}{インストール}

\begin{translateing}
If not available, move inlinelabel.sty file to\\\code{\$TEXMF/tex/latex/inlinelabel}.

直ちに使えなければ，inlinelabel.sty を\\\code{\$TEXMF/tex/latex/inlinelabel}\\%（\TeX が見つけられる場所）
に置いてください。
\end{translateing}

\psuedosection{Loading}{Loading}{読み込み}

\begin{translateing}
To use this package, load .sty file with |\usepackage{inlinelabel}| command in
preamble.

このパッケージを使用するには，プリアンブルに\commandtojskip|\usepackage{inlinelabel}|\commandtojskip と書いてください。
%\end{translateing}

%\begin{translateing}
There are three options:\\
\textbullet\ |[nospace]| swaps with no space command and with space command.\\
\textbullet\ |[circled]| switches to circled equation numbers on \pLaTeX, which is an engine for Japanese.\\
\textbullet\ |[luacircled]| switches to circled equation numbers on \LuaLaTeX\ (and Japanese environment required).

3つのオプションがあります。\\
\textbullet\ |[nospace]|はスペースなし命令とスペースあり命令が入れ替えます。\\
\textbullet\ |[circled]|は丸で囲まれた数式番号に切り替えます（\pLaTeX 用）。\\
\textbullet\ |[luacircled]|は丸で囲まれた数式番号に切り替えます（\LuaLaTeX 用）。
\end{translateing}



\psuedosection{Usage}{Usage}{使用方法}

\macroexplanation{\inlinelabel{<label>}}

\begin{translateing}
This command puts the inline equation number. It takes the name of the label as an argument. Place it outside of the equation (outside of the \$--\$). It makes a small margin before and after the text to make it look natural (in Japanese). If you don't want margins, use |\inlinelabel*{<label>}|.

インラインの数式番号を置きます。引数にラベルの名前を取ります。数式（\$--\$）の外に置いてください。自然に見せるために，前後に少しの余白を取ります。余白がいらないときは\commandtojskip|\inlinelabel*{<label>}|\commandtojskip を使ってください。
%
%If you don't want margins (not in Japanese), use |\inlinelabel*{<label>}|.
%
%日本語組版を念頭に作っていますので，英語では |\inlinelabel*{<label>}| を使って原稿でスペースを入れるのが自然でしょう。

For inconvenience, there is the |[nospace]| option.
In this case, |\inlinelabel{<label>}| is swapped with no space and |\inlinelabel*{<label>}| has space.

それが面倒なときのために，|[nospace]|オプションが用意されています。
このとき，|\inlinelabel{<label>}|\commandtojskip がスペースなし，|\inlinelabel*{<label>}|\commandtojskip がスペースありに入れ替わります。
\end{translateing}

\begin{lstlisting}
An inlinelabel makes 
$e^{i\theta}=\cos\theta+i\sin\theta$\inlinelabel{test1}.

An asterisked one makes 
$e^{i\pi}+1=0$\inlinelabel*{test2}.
\end{lstlisting}

\begin{macroexample}
An inlinelabel makes $e^{i\theta}=\cos\theta+i\sin\theta$\hspace{0.5em}(1)\hspace{0.5em}.

An asterisked one makes $e^{i\pi}+1=0$(2).
\end{macroexample}

\begin{translateing}
The |<label>| can be used for |\ref|.

|<label>|\commandtojskip を\commandtojskip|\ref|で使うことができます。
\end{translateing}

\begin{lstlisting}
Substituting $\theta=\pi$
for (\ref{test1}) yields (\ref{test2}).
\end{lstlisting}

\begin{macroexample}
Substituting $\theta=\pi$ for (\ref*{test1}) yields (\ref*{test2}).
\end{macroexample}

\macroexplanation{Options [circled] [luacircled]}

\begin{translateing}
Assuming you are using a Japanese environment. Switch to circled equation numbers. All equation numbers are rewritten, not just the inline equation numbers.

あなたが日本語の環境を使用していることを前提とします。丸で囲まれた数式番号に切り替えます。インラインの数式番号だけでなく，すべての数式番号が書き換えられます。

When using \pLaTeX\ with the \code{japanese-otf} package, specify \code{[circled]}; when using \LuaLaTeX\ with the \code{luatexja-otf} package, specify \code{[luacircled]}.

\pLaTeX を使っているときは，\code{japanese-otf} パッケージの読み込みとともに\code{[circled]}オプションを指定してください。
\LuaLaTeX を使っているときは，\code{luatexja-otf} パッケージの読み込みとともに\code{[luacircled]}オプションを指定してください。
\end{translateing}
%
%\begin{lstlisting}
%$e^{i\theta}=\cos\theta+i\sin\theta$\inlinelabel{test1}はテストです。
%
%アスタリスクをつけると$e^{i\pi}+1=0$\inlinelabel*{test2}です。
%\end{lstlisting}

\begin{macroexample}
$e^{i\theta}=\cos\theta+i\sin\theta$\inlinelabel{test1}はテストです。

アスタリスクをつけると$e^{i\pi}+1=0$\inlinelabel*{test2}です。
\end{macroexample}

\begin{translateing}
The following example is just an align environment.

次の例は単なる align 環境です。
\end{translateing}

\begin{lstlisting}
\begin{align}
&\phantom{=}\log_{3}x+\log_{3}(6-x)\\
&=\log_{3}x(6-x)\\
&=\log_{3}(-(x-3)^{2}+9)
\end{align}
\end{lstlisting}

%\equationreset
\begin{macroexample*}{align}
&\phantom{=}\log_{3}x+\log_{3}(6-x)\\
&=\log_{3}x(6-x)\\
&=\log_{3}(-(x-3)^{2}+9)
\end{macroexample*}

\macroexplanation{\circledref{<label>}  \equationref{<label>}}

\begin{translateing}
These commands are only valid when \code{[circled]} or \code{[luacircled]} option is loaded.
They refer to the circled equation number.

\code{[circled]}または\code{[luacircled]}オプションを読み込んでいるときにのみ有効です。
丸で囲まれた数式番号を参照します。
\end{translateing}

\begin{translateing}
|\equationref| is synonym with |\circledref|.

|\equationref|は\commandtojskip|\circledref|の別名です。
\end{translateing}

\begin{lstlisting}
Substituting $\theta=\pi$
for \equationref{test1} yields \equationref{test2}.
\end{lstlisting}

\begin{macroexample}
Substituting $\theta=\pi$ for \equationref{test1} yields \equationref{test2}.
\end{macroexample}

\macroexplanation{\equationreset}

\begin{translateing}
Reset the equation number back to 1.

数式番号を1に戻します。
\end{translateing}

\psuedosection{moreinfo}{For More Information}{問い合わせ・詳しくは}

\noindent\hspace*{\fill}\begin{tabular}{rl}
The inlinelabel package:&\url{https://www.metaphysica.info/technote/package_inlinelabel/}\\
Yukoh KUSAKABE:&\url{https://www.metaphysica.info/}\\
&\url{https://twitter.com/metaphysicainfo}\\
&(screen-name, 日下部幽考 in Japanese)
\end{tabular}\hspace*{\fill}
\end{document}